\subsection{Odchylka přímek}
\label{subsec:Angle_between_lines}

V předchozích částech jsme určovali, zda jsou dvě přímky různoběžné,
rovnoběžné, totožné nebo kolmé. U různoběžných přímek však můžeme
zkoumat také velikost úhlu, pod kterým se protínají.

Směr přímky je určen jejím směrovým vektorem. Výpočet odchylky dvou
přímek proto navazuje na geometrický význam skalárního součinu,
kterému jsme se věnovali v části
\ref{subsubsec:Geometric_meaning_of_dot_product}. Budeme přitom
vycházet zejména ze vztahu \eqref{eq:dot_product_angle} pro výpočet
úhlu dvou vektorů, případně z jeho odvozené formy z poznámky \ref{rem:dot_product_angle}


\subsubsection{Geometrický význam odchylky přímek}
\label{subsubsec:Geometric_meaning_angle_lines}

Dvě různoběžné přímky vytvářejí ve svém průsečíku čtyři úhly.
Protilehlé úhly mají stejnou velikost a dva sousední úhly jsou
vedlejší, takže jejich součet je roven $\pi$.

Při určování odchylky přímek vybíráme menší z úhlů, které přímky
svírají. Jsou-li přímky kolmé, je jejich odchylka rovna
$\frac{\pi}{2}$. Rovnoběžné a totožné přímky mají odchylku rovnou
nule.

\begin{definition}\label{def:Angle_between_lines}
    \textbf{Odchylkou dvou přímek} $p$ a $q$ v rovině nazýváme číslo
    \[
        \varphi\in\left[0;\frac{\pi}{2}\right],
    \]
    které je u různoběžných přímek velikostí jejich ostrého nebo
    pravého úhlu.

    Pro rovnoběžné a totožné přímky definujeme
    \[
        \varphi=0.
    \]
\end{definition}

Odchylku přímek budeme značit
\[
    \varphi=\angle(p,q).
\]

\begin{figure}[H]
    \centering
    \begin{tikzpicture}[scale=1.05]

        \coordinate (O) at (0,0);

        % Přímky
        \draw[thin, <->] (-3,-1.5) -- (3,1.5)
            node[above right] {$p$};
        \draw[thin, <->] (-2.6,2.2) -- (2.6,-2.2)
            node[below right] {$q$};

        % Průsečík
        \fill (O) circle (2.3pt) node[below left] {$P$};

        % Směrové vektory
        \draw[vector] (0.4,0.2) -- (1.8,0.9)
            node[midway, above, sloped] {$\vec{u}$};
        \draw[vector] (0.35,-0.30) -- (1.55,-1.31)
            node[midway, below, sloped] {$\vec{v}$};

        % Oblouk odchylky
        \draw[->]
            (0.85,0.425)
            arc[start angle=26.565, end angle=-40.236, radius=0.95];

        \node[right] at (0.92,-0.08) {$\varphi$};

    \end{tikzpicture}
    \caption{Odchylka přímek je velikost menšího úhlu, který přímky
    svírají.}
    \label{fig:Angle_between_lines}
\end{figure}

Úhel dvou vektorů může nabývat hodnot od $0$ do $\pi$. Odchylka
přímek však nepřesahuje $\frac{\pi}{2}$. Navíc můžeme směrový vektor
přímky nahradit jeho opačným vektorem, aniž bychom změnili směr
přímky. Výpočet odchylky proto nesmí záviset na zvolené orientaci
směrových vektorů.


\subsubsection{Odchylka pomocí směrových vektorů}
\label{subsubsec:Angle_lines_direction_vectors}

Uvažujme přímky $p$ a $q$ s nenulovými směrovými vektory
\[
    \vec{u}_p=(u_1;u_2),
    \qquad
    \vec{u}_q=(v_1;v_2).
\]

Pro úhel $\alpha\in[0;\pi]$ jejich směrových vektorů platí
\[
    \cos\alpha
    =
    \frac{\vec{u}_p\cdot\vec{u}_q}
    {|\vec{u}_p|\,|\vec{u}_q|}.
\]

Je-li úhel $\alpha$ ostrý nebo pravý, je přímo odchylkou přímek.
Je-li $\alpha$ tupý, je odchylkou přímek vedlejší úhel
\[
    \varphi=\pi-\alpha.
\]
Pro jejich kosiny platí
\[
    \cos(\pi-\alpha)=-\cos\alpha.
\]

Obě možnosti můžeme spojit pomocí absolutní hodnoty skalárního
součinu. Pro odchylku přímek tedy platí
\begin{equation}\label{eq:Angle_lines_direction_vectors}
        \cos\varphi
    =
    \frac{
        \left|\vec{u}_p\cdot\vec{u}_q\right|
    }
    {
        |\vec{u}_p|\,|\vec{u}_q|
    }.
\end{equation}

Po rozepsání souřadnic směrových vektorů dostáváme
\begin{equation}\label{eq:Angle_lines_direction_vectors_coordinates}
    \cos\varphi
    =
    \frac{
        |u_1v_1+u_2v_2|
    }
    {
        \sqrt{u_1^2+u_2^2}\,
        \sqrt{v_1^2+v_2^2}
    }.
\end{equation}

Samotnou odchylku určíme pomocí funkce arkus kosinus:
\begin{equation}\label{eq:Angle_lines_arccos}
    \varphi
    =
    \arccos
    \left(
        \frac{
            \left|\vec{u}_p\cdot\vec{u}_q\right|
        }
        {
            |\vec{u}_p|\,|\vec{u}_q|
        }
    \right).
\end{equation}

\begin{remark}\label{rem:Absolute_value_angle_lines}
    Absolutní hodnota v čitateli zajišťuje dvě důležité vlastnosti:

    \begin{itemize}
        \item výsledná odchylka nepřesáhne $\frac{\pi}{2}$;
        \item výsledek nezávisí na orientaci směrových vektorů.
    \end{itemize}

    Nahradíme-li například vektor $\vec{u}_p$ opačným vektorem
    $-\vec{u}_p$, změní se znaménko skalárního součinu, ale jeho
    absolutní hodnota zůstane stejná.
\end{remark}

Ze vztahu \eqref{eq:Angle_lines_direction_vectors} přímo plynou známé
zvláštní případy:

\[
\begin{array}{c|c|c}
    \text{podmínka}
    &
    \cos\varphi
    &
    \text{vzájemná poloha}
    \\
    \hline
    \vec{u}_p=k\vec{u}_q,\ k\neq0
    &
    1
    &
    \varphi=0
    \\[1mm]
    \vec{u}_p\cdot\vec{u}_q=0
    &
    0
    &
    \varphi=\frac{\pi}{2}
\end{array}
\]

Rovnoběžné a totožné přímky mají odchylku $0$, zatímco kolmé přímky
mají odchylku $\frac{\pi}{2}$.

\begin{example}\label{exa:Angle_parametric_lines}
    Určete odchylku přímek
    \[
        \begin{aligned}
            p:\quad x&=1+2t,\\
                     y&=3-3t,
                     \qquad t\in\mathbb{R},
        \end{aligned}
    \]
    a
    \[
        \begin{aligned}
            q:\quad x&=2-s,\\
                     y&=3+s,
                     \qquad s\in\mathbb{R}.
        \end{aligned}
    \]

    Z parametrických vyjádření získáme směrové vektory
    \[
        \vec{u}_p=(2;-3),
        \qquad
        \vec{u}_q=(-1;1).
    \]

    Jejich skalární součin je
    \[
        \vec{u}_p\cdot\vec{u}_q
        =
        (2;-3)\cdot(-1;1)
        =
        -2-3
        =
        -5.
    \]

    Velikosti směrových vektorů jsou
    \[
        |\vec{u}_p|
        =
        \sqrt{2^2+(-3)^2}
        =
        \sqrt{13}
    \]
    a
    \[
        |\vec{u}_q|
        =
        \sqrt{(-1)^2+1^2}
        =
        \sqrt2.
    \]

    Dosadíme do vztahu
    \eqref{eq:Angle_lines_direction_vectors}:
    \begin{align*}
        \cos\varphi
        &=
        \frac{|-5|}{\sqrt{13}\sqrt2}\\
        &=
        \frac{5}{\sqrt{26}}.
    \end{align*}

    Odtud
    \[
        \varphi
        =
        \arccos\left(\frac{5}{\sqrt{26}}\right)
        \doteq
        11{,}31^\circ.
    \]

    Odchylka přímek je přibližně
    \[
        \varphi\doteq11^\circ19'.
    \]
\end{example}


\subsubsection{Odchylka přímek zadaných obecnými rovnicemi}
\label{subsubsec:Angle_general_lines}

Uvažujme přímky zadané obecnými rovnicemi
\[
    p:\quad a_1x+b_1y+c_1=0
\]
a
\[
    q:\quad a_2x+b_2y+c_2=0.
\]

Jejich normálové vektory jsou
\[
    \vec{n}_p=(a_1;b_1),
    \qquad
    \vec{n}_q=(a_2;b_2).
\]

Obě přímky otočíme o pravý úhel, zvolíme-li místo nich přímky ve
směru jejich normálových vektorů. Otočení obou směrů o stejný úhel
nezmění jejich vzájemnou odchylku. Odchylka přímek $p$ a $q$ je proto
stejná jako odchylka jejich normálových vektorů.

Můžeme tedy použít vztah
\begin{equation}\label{eq:Angle_general_lines}
    \cos\varphi
    =
    \frac{
        |a_1a_2+b_1b_2|
    }
    {
        \sqrt{a_1^2+b_1^2}\,
        \sqrt{a_2^2+b_2^2}
    }.
\end{equation}

Absolutní členy $c_1$ a $c_2$ se ve vzorci nevyskytují, protože
ovlivňují polohu přímek, nikoli jejich směr.

\begin{figure}[H]
    \centering
    \begin{tikzpicture}[scale=1.0]

        \coordinate (O) at (0,0);

        % Přímky
        \draw[thin, <->] (-3,-1.25) -- (3,1.25)
            node[above right] {$p$};
        \draw[thin, <->] (-2.7,2.0) -- (2.7,-2.0)
            node[below right] {$q$};

        % Normálové vektory z průsečíku
        \draw[vector] (O) -- (-0.7,1.7)
            node[midway, left] {$\vec{n}_p$};
        \draw[vector] (O) -- (1.25,1.7)
            node[midway, right] {$\vec{n}_q$};

        % Oblouky
        \draw[->]
            (0.75,0.313)
            arc[start angle=22.62, end angle=-36.53, radius=0.81];
        \node[right] at (0.82,-0.05) {$\varphi$};

        \draw[->]
            (-0.32,0.78)
            arc[start angle=112.38, end angle=53.67, radius=0.84];
        \node[above] at (0.05,0.86) {$\varphi$};

        \fill (O) circle (2.2pt);

    \end{tikzpicture}
    \caption{Odchylka dvou přímek je stejná jako odchylka jejich
    normálových vektorů.}
    \label{fig:Angle_general_lines_normals}
\end{figure}

\begin{example}\label{exa:Angle_general_line_equations}
    Určete odchylku přímek $p: 2x-y+3=0$ a $q:3x+2y-1=0$.

    Normálové vektory přímek jsou
    \[
        \vec{n}_p=(2;-1),
        \qquad
        \vec{n}_q=(3;2).
    \]

    Jejich skalární součin je
    \[
        \vec{n}_p\cdot\vec{n}_q
        =
        (2;-1)\cdot(3;2)
        =
        6-2
        =
        4.
    \]

    Velikosti normálových vektorů jsou
    \[
        |\vec{n}_p|=\sqrt{2^2+(-1)^2}=\sqrt5
    \]
    a
    \[
        |\vec{n}_q|=\sqrt{3^2+2^2}=\sqrt{13}.
    \]

    Dosadíme do vztahu \eqref{eq:Angle_general_lines}:
    \[
        \cos\varphi
        =
        \frac{4}{\sqrt5\sqrt{13}}
        =
        \frac{4}{\sqrt{65}}.
    \]

    Odtud
    \[
        \varphi
        =
        \arccos\left(\frac{4}{\sqrt{65}}\right)
        \doteq
        60{,}26^\circ.
    \]

    Odchylka přímek je přibližně
    \[
        \varphi\doteq60^\circ15'.
    \]
\end{example}


\subsubsection{Odchylka přímek ve směrnicovém tvaru}
\label{subsubsec:Angle_slope_lines}

Uvažujme dvě nesvislé přímky ve směrnicovém tvaru
\[
    p:\quad y=k_1x+q_1
\]
a
\[
    q:\quad y=k_2x+q_2.
\]

Jako jejich směrové vektory můžeme zvolit
\[
    \vec{u}_p=(1;k_1),
    \qquad
    \vec{u}_q=(1;k_2).
\]

Dosazením do vztahu
\eqref{eq:Angle_lines_direction_vectors} získáme
\begin{equation}\label{eq:Angle_slope_lines_cosine}
    \cos\varphi
    =
    \frac{|1+k_1k_2|}
    {\sqrt{1+k_1^2}\sqrt{1+k_2^2}}.
\end{equation}

Odchylku můžeme určit také pomocí směrových úhlů přímek. Platí-li
\[
    k_1=\tan\alpha_1
    \qquad\text{a}\qquad
    k_2=\tan\alpha_2,
\]
pak s využitím vztahu pro tangens rozdílu dostáváme
\begin{equation}\label{eq:Angle_slope_lines_tangent}
    \tan\varphi
    =
    \left|
        \frac{k_2-k_1}{1+k_1k_2}
    \right|.
\end{equation}

Tento vztah můžeme použít, pokud
\[
    1+k_1k_2\neq0.
\]

Jestliže
\[
    1+k_1k_2=0,
\]
platí
\[
    k_1k_2=-1,
\]
a přímky jsou kolmé. Jejich odchylka je tedy
\[
    \varphi=\frac{\pi}{2}.
\]

Absolutní členy $q_1$ a $q_2$ opět nemají na odchylku vliv.

\begin{remark}\label{rem:Alternative_method_angle_slope_lines}
    Vztahy \eqref{eq:Angle_slope_lines_cosine}
    a \eqref{eq:Angle_slope_lines_tangent} umožňují určit odchylku
    přímo ze směrnic obou přímek.

    Tyto vztahy však není nutné pamatovat jako další samostatné
    vzorce. Obě přímky můžeme nejprve převést ze směrnicového tvaru
    na obecné rovnice a následně použít již známý vztah
    \eqref{eq:Angle_general_lines} pro výpočet odchylky pomocí
    normálových vektorů.

    Přímku
    \[
        y=kx+q
    \]
    převedeme na obecnou rovnici ve tvaru
    \[
        kx-y+q=0.
    \]
    Jejím normálovým vektorem je potom
    \[
        \vec{n}=(k;-1).
    \]

    Oba postupy vedou ke stejnému výsledku. Přímý vztah využívající
    směrnice představuje pouze zkrácení výpočtu.
\end{remark}

\begin{example}\label{exa:Angle_slope_intercept_lines}
    Určete odchylku přímek
    \[
        p:\quad y=2x-1
    \]
    a
    \[
        q:\quad y=\frac12x+3.
    \]

    Směrnice přímek jsou
    \[
        k_1=2,
        \qquad
        k_2=\frac12.
    \]

    Použijeme vztah
    \eqref{eq:Angle_slope_lines_tangent}:
    \begin{align*}
        \tan\varphi
        &=
        \left|
            \frac{\frac12-2}
            {1+2\cdot\frac12}
        \right|\\
        &=
        \left|
            \frac{-\frac32}{2}
        \right|\\
        &=
        \frac34.
    \end{align*}

    Odtud
    \[
        \varphi
        =
        \arctan\left(\frac34\right)
        \doteq
        36{,}87^\circ.
    \]

    Odchylka přímek je přibližně
    \[
        \varphi\doteq36^\circ52'.
    \]
\end{example}

\begin{remark}\label{rem:Vertical_line_angle}
    Směrnicový tvar nelze použít pro svislou přímku. Je-li alespoň
    jedna z přímek svislá, určíme odchylku pomocí směrových nebo
    normálových vektorů. Pro svislou přímku můžeme zvolit například
    směrový vektor
    \[
        \vec{u}=(0;1).
    \]
\end{remark}

\begin{example}\label{exa:Angle_slope_lines_by_conversion}
    Určete odchylku přímek $p: y=\frac{3}{4}x-2$ a $q: y=-\frac{1}{2}x+1$.

    Přímky jsou zadány ve směrnicovém tvaru. Místo přímého použití
    vztahu pro jejich směrnice je nejprve převedeme na obecné rovnice
    a následně použijeme vztah \eqref{eq:Angle_general_lines}.

    Nejprve upravíme rovnici přímky $p$:
    \begin{align*}
        y&=\frac{3}{4}x-2,\\
        \frac{3}{4}x-y-2&=0.
    \end{align*}

    Abychom odstranili zlomek, vynásobíme rovnici číslem $4$:
    \[
        p:\quad 3x-4y-8=0.
    \]

    Normálový vektor přímky $p$ je tedy
    \[
        \vec{n}_p=(3;-4).
    \]

    Nyní převedeme na obecnou rovnici přímku $q$:
    \begin{align*}
        y&=-\frac{1}{2}x+1,\\
        \frac{1}{2}x+y-1&=0.
    \end{align*}

    Po vynásobení rovnice číslem $2$ dostáváme
    \[
        q:\quad x+2y-2=0.
    \]

    Normálový vektor přímky $q$ je
    \[
        \vec{n}_q=(1;2).
    \]

    Vypočítáme skalární součin normálových vektorů:
    \[
        \vec{n}_p\cdot\vec{n}_q
        =
        (3;-4)\cdot(1;2)
        =
        3-8
        =
        -5.
    \]

    Velikosti normálových vektorů jsou
    \[
        |\vec{n}_p|
        =
        \sqrt{3^2+(-4)^2}
        =
        5
    \]
    a
    \[
        |\vec{n}_q|
        =
        \sqrt{1^2+2^2}
        =
        \sqrt5.
    \]

    Dosadíme do vztahu \eqref{eq:Angle_general_lines}:
    \begin{align*}
        \cos\varphi
        &=
        \frac{
            |\vec{n}_p\cdot\vec{n}_q|
        }
        {
            |\vec{n}_p|\,|\vec{n}_q|
        }\\
        &=
        \frac{|-5|}{5\sqrt5}\\
        &=
        \frac{1}{\sqrt5}.
    \end{align*}

    Odtud
    \[
        \varphi
        =
        \arccos\left(\frac{1}{\sqrt5}\right)
        \doteq
        63{,}43^\circ.
    \]

    Odchylka přímek je přibližně
    \[
        \varphi\doteq63^\circ26'.
    \]

    Ke stejnému výsledku bychom dospěli také přímým použitím vztahu
    \eqref{eq:Angle_slope_lines_tangent}. Převod na obecné rovnice
    nám však umožnil použít již známý postup založený na normálových
    vektorech.
\end{example}


\subsubsection{Přímky zadané různými způsoby}
\label{subsubsec:Angle_lines_mixed_forms}

Jsou-li přímky zadány různými způsoby, není nutné vždy převádět celé
jejich rovnice do stejného tvaru. Často stačí z každého zápisu získat
směrový nebo normálový vektor a následně použít vztah
\eqref{eq:Angle_lines_direction_vectors}.

Můžeme postupovat například takto:

\[
\begin{array}{c|c}
    \text{zadání přímky}
    &
    \text{vhodný vektor}
    \\
    \hline
    \text{dva body }A,B
    &
    \overrightarrow{AB}=B-A
    \\[1mm]
    \text{parametrické vyjádření}
    &
    \vec{u}=(u_x;u_y)
    \\[1mm]
    \text{obecná rovnice }ax+by+c=0
    &
    \vec{n}=(a;b)
    \text{ nebo }
    \vec{u}=(-b;a)
    \\[1mm]
    \text{směrnicový tvar }y=kx+q
    &
    \vec{u}=(1;k)
\end{array}
\]

Pro výpočet musí být oba zvolené vektory stejného typu: buď dva
směrové, nebo dva normálové vektory.

\begin{example}\label{exa:Angle_line_through_points_and_general_line}
    Jsou dány body
    \[
        A[-3;1],
        \qquad
        B[1;2]
    \]
    a přímka
    \[
        p:\quad 2x-y+3=0.
    \]

    Určete odchylku přímky $AB$ a přímky $p$.

    Směrový vektor přímky $AB$ je
    \[
        \vec{u}_{AB}
        =
        \overrightarrow{AB}
        =
        (1-(-3);2-1)
        =
        (4;1).
    \]

    Z obecné rovnice přímky $p$ získáme normálový vektor
    \[
        \vec{n}_p=(2;-1).
    \]
    K němu zvolíme kolmý směrový vektor
    \[
        \vec{u}_p=(1;2).
    \]

    Skalární součin směrových vektorů je
    \[
        \vec{u}_{AB}\cdot\vec{u}_p
        =
        (4;1)\cdot(1;2)
        =
        4+2
        =
        6.
    \]

    Jejich velikosti jsou
    \[
        |\vec{u}_{AB}|=\sqrt{4^2+1^2}=\sqrt{17}
    \]
    a
    \[
        |\vec{u}_p|=\sqrt{1^2+2^2}=\sqrt5.
    \]

    Proto
    \[
        \cos\varphi
        =
        \frac{6}{\sqrt{17}\sqrt5}
        =
        \frac{6}{\sqrt{85}}.
    \]

    Odtud
    \[
        \varphi
        =
        \arccos\left(\frac{6}{\sqrt{85}}\right)
        \doteq
        49{,}40^\circ.
    \]

    Odchylka přímek je přibližně
    \[
        \varphi\doteq49^\circ24'.
    \]
\end{example}





\subsubsection{Přímka procházející bodem s danou odchylkou}
\label{subsubsec:Line_through_point_given_angle}

Úlohu můžeme také obrátit. Známe-li přímku $p$, bod $Q$ a požadovanou
odchylku $\varphi$, hledáme přímky procházející bodem $Q$, které
svírají s přímkou $p$ daný úhel.

Není-li požadovaná odchylka rovna nule ani pravému úhlu, existují
obvykle dvě takové přímky. Jedna vznikne otočením směru přímky $p$
o úhel $\varphi$ na jednu stranu a druhá otočením o tentýž úhel na
stranu opačnou.

\begin{example}\label{exa:Lines_through_point_given_angle}
    Je dána přímka $p:x-3y-2=0$ a bod $Q[1;1]$. Určete všechny
    přímky procházející bodem $Q$, jejichž odchylka od přímky $p$
    je $45^\circ$.

    Z obecné rovnice přímky $p$ získáme normálový vektor
    \[
        \vec{n}_p=(1;-3).
    \]
    Jako směrový vektor přímky $p$ můžeme zvolit
    \[
        \vec{u}_p=(3;1).
    \]

    Hledaná přímka $q$ prochází bodem $Q[1;1]$. Nejprve budeme
    předpokládat, že není svislá, a její směrový vektor zapíšeme jako
    \[
        \vec{u}_q=(1;k),
    \]
    kde $k$ je její směrnice.

    Pro odchylku $45^\circ$ musí platit
    \[
        \cos45^\circ
        =
        \frac{
            |\vec{u}_p\cdot\vec{u}_q|
        }
        {
            |\vec{u}_p|\,|\vec{u}_q|
        }.
    \]

    Po dosazení dostáváme
    \[
        \frac{\sqrt2}{2}
        =
        \frac{|3+k|}
        {\sqrt{10}\sqrt{1+k^2}}.
    \]

    Umocníme obě strany rovnice a postupně upravíme:
    \begin{align*}
        \frac12
        &=
        \frac{(3+k)^2}{10(1+k^2)},\\
        5(1+k^2)
        &=
        (3+k)^2,\\
        5+5k^2
        &=
        9+6k+k^2,\\
        4k^2-6k-4
        &=
        0,\\
        2k^2-3k-2
        &=
        0.
    \end{align*}

    Vyřešíme kvadratickou rovnici:
    \[
        (2k+1)(k-2)=0.
    \]

    Dostáváme dvě směrnice:
    \[
        k_1=2,
        \qquad
        k_2=-\frac12.
    \]

    \paragraph{První hledaná přímka}

    Přímka se směrnicí $k_1=2$ prochází bodem $Q[1;1]$. Použijeme
    bodově-směrnicový tvar (viz vztah \eqref{eq:Point_slope_line_equation}):
    \begin{align*}
        y-1&=2(x-1),\\
        y-1&=2x-2,\\
        2x-y-1&=0.
    \end{align*}

    První hledaná přímka je
    \[
        q_1:2x-y-1=0.
    \]

    Přímka $q_1$ společně s původní přímkou $p$ je znázorněna
    na obrázku~\ref{fig:First_line_through_point_given_angle}.
    Přímky se protínají v bodě
    \[
        P_1\left[\frac15;-\frac35\right]
    \]
    a jejich odchylka je $45^\circ$.

    \begin{figure}[H]
        \centering
        \begin{tikzpicture}[scale=0.85]

            % Osy
            \draw[<->, thin] (-3,0) -- (6,0) node[right] {$x$};
            \draw[<->, thin] (0,-3) -- (0,6) node[above] {$y$};
            \node[below left] at (0,0) {$O$};

            % Původní přímka p: y=(x-2)/3
            \draw[thin, <->]
                (-3,-5/3) -- (6,4/3)
                node[above right] {$p$};

            % První hledaná přímka q_1: y=2x-1
            \draw[thick, <->]
                (-1,-3) -- (3.5,6)
                node[above right] {$q_1$};

            % Bod Q
            \fill (1,1) circle (2.6pt)
                node[above left] {$Q[1;1]$};

            % Průsečík p a q_1
            \coordinate (Pone) at (0.2,-0.6);
            \fill (Pone) circle (2.2pt)
                node[below right] {$P_1$};

            % Odchylka 45 stupňů
            \draw[->]
                ($(Pone)+(0.75,0.25)$)
                arc[
                    start angle=18.435,
                    end angle=63.435,
                    radius=0.79
                ];

            \node[right] at (0.85,0.02) {$45^\circ$};

        \end{tikzpicture}
        \caption{První přímka $q_1$ procházející bodem $Q$ a svírající
        s přímkou $p$ úhel $45^\circ$.}
        \label{fig:First_line_through_point_given_angle}
    \end{figure}


    \paragraph{Druhá hledaná přímka}

    Pro směrnici $k_2=-\frac12$ dostáváme
    \begin{align*}
        y-1&=-\frac12(x-1),\\
        2y-2&=-x+1,\\
        x+2y-3&=0.
    \end{align*}

    Druhá hledaná přímka je
    \[
        q_2:x+2y-3=0.
    \]

    Na obrázku~\ref{fig:Second_line_through_point_given_angle} je
    nově sestrojená přímka $q_2$ znázorněna plnou čarou. První
    nalezená přímka $q_1$ je pro porovnání vyznačena čárkovaně.
    Přímky $p$ a $q_2$ se protínají v bodě
    \[
        P_2\left[\frac{13}{5};\frac15\right]
    \]
    a jejich odchylka je opět $45^\circ$.

    \begin{figure}[H]
        \centering
        \begin{tikzpicture}[scale=0.85]

            % Osy
            \draw[<->, thin] (-3,0) -- (6,0) node[right] {$x$};
            \draw[<->, thin] (0,-3) -- (0,6) node[above] {$y$};
            \node[below left] at (0,0) {$O$};

            % Původní přímka p
            \draw[thin, <->]
                (-3,-5/3) -- (6,4/3)
                node[above right] {$p$};

            % První řešení čárkovaně
            \draw[thick, dashed, <->]
                (-1,-3) -- (3.5,6)
                node[above right] {$q_1$};

            % Druhé řešení plnou čarou
            \draw[thick, <->]
                (-3,3) -- (6,-1.5)
                node[below right] {$q_2$};

            % Bod Q
            \fill (1,1) circle (2.6pt)
                node[above right] {$Q[1;1]$};

            % Průsečík p a q_2
            \coordinate (Ptwo) at (2.6,0.2);
            \fill (Ptwo) circle (2.2pt)
                node[above right] {$P_2$};

            % Odchylka 45 stupňů
            \draw[->]
                ($(Ptwo)+(0.72,-0.36)$)
                arc[
                    start angle=-26.565,
                    end angle=18.435,
                    radius=0.81
                ];

            \node[right] at (3.38,0.18) {$45^\circ$};

        \end{tikzpicture}
        \caption{Druhá přímka $q_2$ procházející bodem $Q$. První
        nalezená přímka $q_1$ je znázorněna čárkovaně.}
        \label{fig:Second_line_through_point_given_angle}
    \end{figure}

    Svislá přímka procházející bodem $Q$ má směrový vektor
    \[
        \vec{u}_s=(0;1).
    \]
    Pro její odchylku od přímky $p$ by platilo
    \[
        \cos\varphi
        =
        \frac{|(3;1)\cdot(0;1)|}
        {\sqrt{3^2+1^2}\sqrt{0^2+1^2}}
        =
        \frac{1}{\sqrt{10}}.
    \]
    Tato hodnota neodpovídá úhlu $45^\circ$, takže žádné další
    řešení nevzniká.

    Zadání splňují právě dvě přímky:
    \[
        q_1:2x-y-1=0
    \]
    a
    \[
        q_2:x+2y-3=0.
    \]

    Celková situace je znázorněna na obrázku
    \ref{fig:Lines_through_point_given_angle}. Původní přímka $p$
    je v obrázku zvýrazněna silnější čarou.

    \begin{figure}[H]
        \centering
        \begin{tikzpicture}[scale=1]

            % Osy
            \draw[<->, thin] (-4,0) -- (7,0) node[right] {$x$};
            \draw[<->, thin] (0,-4) -- (0,7) node[above] {$y$};
            \node[below left] at (0,0) {$O$};

            % Hodnoty na ose x
            \foreach \x in {-3,-2,-1,1,2,3,4,5,6}
            {
                \draw (\x,0.10) -- (\x,-0.10);
                \node[below, font=\scriptsize] at (\x,0) {$\x$};
            }

            % Hodnoty na ose y
            \foreach \y in {-3,-2,-1,1,2,3,4,5,6}
            {
                \draw (0.10,\y) -- (-0.10,\y);
                \node[left, font=\scriptsize] at (0,\y) {$\y$};
            }

            % Původní přímka p zvýrazněně
            \draw[very thick, <->]
                (-4,-2) -- (7,5/3)
                node[above right] {$p$};

            % Hledané přímky
            \draw[thick, <->]
                (-1.5,-4) -- (4,7)
                node[above right] {$q_1$};

            \draw[thick, <->]
                (-4,3.5) -- (7,-2)
                node[below right] {$q_2$};

            % Bod Q
            \fill (1,1) circle (3pt)
                node[above right] {$Q[1;1]$};

            % Průsečíky
            \coordinate (Pone) at (0.2,-0.6);
            \coordinate (Ptwo) at (2.6,0.2);

            \fill (Pone) circle (2.2pt)
                node[below left] {$P_1$};

            \fill (Ptwo) circle (2.2pt)
                node[above right] {$P_2$};

            % První úhel
            \draw[->]
                ($(Pone)+(0.72,0.24)$)
                arc[
                    start angle=18.435,
                    end angle=63.435,
                    radius=0.76
                ];

            \node[above right, font=\small] at (0.82,0.00) {$45^\circ$};

            % Druhý úhel
            \draw[->]
                ($(Ptwo)+(0.68,-0.34)$)
                arc[
                    start angle=-26.565,
                    end angle=18.435,
                    radius=0.76
                ];

            \node[right, font=\small] at (3.34,0.18) {$45^\circ$};

        \end{tikzpicture}
        \caption{Přímky $q_1$ a $q_2$ procházející bodem $Q$, jejichž
        odchylka od zvýrazněné přímky $p$ je $45^\circ$.}
        \label{fig:Lines_through_point_given_angle}
    \end{figure}
\end{example}


\begin{remark}\label{rem:Summary_angle_between_lines}
    Při výpočtu odchylky přímek není nutné pamatovat několik
    nesouvisejících vzorců. Základním postupem je vždy určit dva
    směrové vektory a použít vztah
    \eqref{eq:Angle_lines_direction_vectors}.

    Vzorce pro obecné rovnice a směrnicové tvary jsou pouze jeho
    upravenými podobami, které umožňují pracovat přímo se zadaným
    tvarem přímek.
\end{remark}
